%% kapa: exam=1 round=29 year=2018 period=1 form=A subject=economics subject_ko=경제학원론 seq=15 number=55
%% source: Q-Net 감정평가사 기출문제 (cst00309) artlSeq=5202808 file=extracted/1008/5202808/2018년 제29회 감정평가사 1차시험 기출문제/2018년 제29회 감정평가사 1차시험 기출문제/1교시 A형.pdf
%% conversion: kordoc 4.24.51 -> tools/gen1.py (auto); review=proofread-sonnet; flags=none
\begin{question}{55}
효용을 극대화하는 甲은 1기의 소비($c_{1}$)와 2기의 소비( $c_{2}$)로 구성된 효용함수
$U(c_{1},c_{2})=c_{1}c_{2}^{2}$을 가지고 있다. 甲은 시점 간 선택(intertemporal choice) 모형에서 1기에 3,000만 원, 2기에 3,300만 원의 소득을 얻고, 이자율 10 \%로 저축하거나 빌릴 수 있다. 1기의 최적 선택에 관한 설명으로 옳은 것은? (단, 인플레이션은 고려하지 않는다.) \par
\begin{choices}
  \item 1,000만 원을 저축할 것이다.
  \item 1,000만 원을 빌릴 것이다.
  \item 저축하지도 빌리지도 않을 것이다.
  \item 1,400만 원을 저축할 것이다.
  \item 1,400만 원을 빌릴 것이다.
\end{choices}
\end{question}
\answer{1}
